\section{分类训练与限时训练}

\subsection{基础巩固训练}

\subsubsection{物质的量 $N_A$ 专题}

\textbf{1.} 设 $N_A$ 为阿伏加德罗常数的值。下列说法正确的是
\begin{enumerate}
  \item[A.] $1\ \text{L}\ 0.1\ \text{mol/L}\ \text{Na}_2\text{CO}_3$ 溶液中含有的 $\text{CO}_3^{2-}$ 数目为 $0.1N_A$
  \item[B.] 标准状况下，$2.24\ \text{L}\ \text{SO}_3$ 含有的氧原子数目为 $0.3N_A$
  \item[C.] $1\ \text{mol}\ \text{Cl}_2$ 与足量 $\text{Fe}$ 反应，转移电子数目为 $2N_A$
  \item[D.] $1\ \text{mol}\ \text{SiO}_2$ 中含有 $\text{Si-O}$ 键的数目为 $2N_A$
\end{enumerate}

\textbf{2.} $N_A$ 为阿伏加德罗常数的值。$2\ \text{mol}\ \text{SO}_2$ 和 $1\ \text{mol}\ \text{O}_2$ 在密闭容器中充分反应，生成 $\text{SO}_3$ 的分子数
A. 等于 $2N_A$ \quad B. 小于 $2N_A$ \quad C. 大于 $2N_A$ \quad D. 无法确定

\subsubsection{离子方程式专题}

\textbf{3.} 下列离子方程式书写正确的是
A. $\text{Cl}_2$ 通入水中：$\text{Cl}_2 + \text{H}_2\text{O} = 2\text{H}^+ + \text{Cl}^- + \text{ClO}^-$
B. 向 $\text{Fe(OH)}_3$ 中加入氢碘酸：$\text{Fe(OH)}_3 + 3\text{H}^+ = \text{Fe}^{3+} + 3\text{H}_2\text{O}$
C. 向 $\text{Ca(ClO)}_2$ 溶液中通入少量 $\text{SO}_2$：$\text{Ca}^{2+} + 2\text{ClO}^- + \text{SO}_2 + \text{H}_2\text{O} = \text{CaSO}_3\downarrow + 2\text{HClO}$
D. $\text{NaHCO}_3$ 溶液与足量 $\text{Ba(OH)}_2$ 反应：$\text{HCO}_3^- + \text{OH}^- + \text{Ba}^{2+} = \text{BaCO}_3\downarrow + \text{H}_2\text{O}$

\subsubsection{元素周期律专题}

\textbf{4.} 短周期主族元素 $X$、$Y$、$Z$、$W$ 原子序数依次增大，$X$ 的简单氢化物和 $Y$ 的氢化物相遇产生白烟，$Z$ 的单质为淡黄色固体，$W$ 的最高价氧化物对应水化物是强酸。下列说法正确的是
A. 原子半径：$X < Y < Z < W$
B. 简单离子半径：$Y < Z < W$
C. $X$ 与 $Z$ 形成的化合物中只含离子键
D. $Y$ 的简单氢化物比 $X$ 的稳定

\subsubsection{有机化学基础}

\textbf{5.} 下列有关有机物的说法正确的是
A. 乙烯和聚乙烯都能使 $\text{Br}_2$ 的 $\text{CCl}_4$ 溶液褪色
B. 乙醛和葡萄糖都能发生银镜反应
C. 苯和苯酚都能与 $\text{Br}_2$ 水发生取代反应
D. 乙醇和乙酸乙酯都能发生消去反应

\textbf{6.} 分子式为 $\text{C}_5\text{H}_{10}\text{O}_2$ 的有机物在酸性条件下可水解为酸和醇，若不考虑立体异构，这些醇和酸重新组合可形成的酯有
A. 20 种 \quad B. 24 种 \quad C. 28 种 \quad D. 32 种

\subsubsection{化学实验基础}

\textbf{7.} 下列实验操作能达到实验目的的是
A. 用酒精萃取碘水中的碘
B. 用蒸馏法分离乙酸乙酯和乙醇
C. 用饱和 $\text{Na}_2\text{CO}_3$ 溶液除去乙酸乙酯中的乙酸
D. 用分液漏斗分离 $\text{CCl}_4$ 和 $\text{I}_2$

\subsection{能力提升训练}

\subsubsection{工业流程综合}

\textbf{8.} 以菱铁矿（主要成分 $\text{FeCO}_3$，含少量 $\text{SiO}_2$）为原料制备高纯 $\text{Fe}_2\text{O}_3$ 的流程如下：

菱铁矿 $\xrightarrow{\text{稀H}_2\text{SO}_4}$ 酸浸 $\to$ 过滤 $\xrightarrow{\text{H}_2\text{O}_2}$ 氧化 $\xrightarrow{\text{调pH}}$ 过滤 $\xrightarrow{\text{煅烧}}$ $\text{Fe}_2\text{O}_3$

\begin{enumerate}
  \item 酸浸时发生反应的离子方程式：\underline{\qquad\qquad}。
  \item 加入 $\text{H}_2\text{O}_2$ 的作用是\underline{\qquad\qquad}，离子方程式：\underline{\qquad\qquad}。
  \item 调 pH 的目的是\underline{\qquad\qquad}。
  \item 煅烧过程的化学方程式：\underline{\qquad\qquad}。
\end{enumerate}

\subsubsection{化学反应原理综合}

\textbf{9.} 甲醇是一种重要的化工原料。工业上利用 $\text{CO}_2$ 和 $\text{H}_2$ 合成甲醇：
\[
\text{CO}_2(g) + 3\text{H}_2(g) \rightleftharpoons \text{CH}_3\text{OH}(g) + \text{H}_2\text{O}(g)\ \Delta H < 0
\]

\begin{enumerate}
  \item 该反应的平衡常数表达式 $K =$\underline{\qquad\qquad}。
  \item 下列措施能提高 $\text{CH}_3\text{OH}$ 产率的是\underline{\qquad\qquad}（填序号）。
  \begin{enumerate}
    \item[A.] 升高温度 \quad B. 增大压强 \quad C. 使用催化剂 \quad D. 增大 $\text{CO}_2$ 浓度
  \end{enumerate}
  \item 某温度下，向 $2\ \text{L}$ 恒容密闭容器中充入 $1\ \text{mol}\ \text{CO}_2$ 和 $3\ \text{mol}\ \text{H}_2$，测得 $\text{CH}_3\text{OH}$ 的物质的量随时间变化如下：
  \[
  \begin{array}{c|cccccc}
  t/\text{min} & 0 & 2 & 4 & 6 & 8 & 10 \\ \hline
  n(\text{CH}_3\text{OH})/\text{mol} & 0 & 0.2 & 0.35 & 0.45 & 0.50 & 0.50
  \end{array}
  \]
  \begin{enumerate}
    \item 反应在第\underline{\qquad\qquad} $\text{min}$ 达到平衡。
    \item $0\text{--}2\ \text{min}$ 内，$v(\text{H}_2) =$\underline{\qquad\qquad} $\text{mol/(L}\cdot\text{min)}$。
    \item 该温度下的平衡常数 $K =$\underline{\qquad\qquad}（保留两位有效数字）。
  \end{enumerate}
\end{enumerate}

\subsubsection{实验探究综合}

\textbf{10.} 某小组设计实验探究 $\text{SO}_2$ 与 $\text{BaCl}_2$ 溶液反应的条件。

\begin{enumerate}
  \item 提出假设：$\text{SO}_2$ 与 $\text{BaCl}_2$ 溶液在\underline{\qquad\qquad}条件下能反应生成白色沉淀。
  \item 设计实验验证假设：
  \[
  \begin{array}{c|c|c}
  \text{实验序号} & \text{实验操作} & \text{预期现象} \\ \hline
  ① & \text{向}\ \text{BaCl}_2\ \text{溶液中通入}\ \text{SO}_2 & \text{无沉淀} \\
  ② & \text{向}\ \text{BaCl}_2\ \text{溶液中通入}\ \text{SO}_2\ \text{和}\ \text{NH}_3 & \underline{\qquad\qquad} \\
  ③ & \text{向}\ \text{BaCl}_2\ \text{溶液中通入}\ \text{SO}_2\ \text{后加入}\ \text{H}_2\text{O}_2 & \underline{\qquad\qquad} \\
  \end{array}
  \]
  \item 实验②中反应的离子方程式：\underline{\qquad\qquad}。
  \item 实验③中反应的离子方程式：\underline{\qquad\qquad}。
\end{enumerate}

\subsubsection{有机推断综合}

\textbf{11.} 已知 $A$ 是芳香族化合物，其分子式为 $\text{C}_7\text{H}_8\text{O}$，$A$ 能发生如下转化：

$A(\text{C}_7\text{H}_8\text{O}) \xrightarrow{\text{浓HBr}} B(\text{C}_7\text{H}_7\text{Br}) \xrightarrow{\text{NaOH/H}_2\text{O}} C \xrightarrow{\text{Cu},\Delta} D \xrightarrow{\text{银氨溶液}} E \xrightarrow{\text{H}^+} F$

\begin{enumerate}
  \item $A$ 的结构简式：\underline{\qquad\qquad}，名称：\underline{\qquad\qquad}。
  \item $A \to B$ 的反应类型：\underline{\qquad\qquad}。
  \item $C \to D$ 的化学方程式：\underline{\qquad\qquad}。
  \item $F$ 与 $\text{NaHCO}_3$ 溶液反应的离子方程式：\underline{\qquad\qquad}。
  \item $F$ 的同分异构体中，属于酯类的芳香族化合物有\underline{\qquad\qquad}种。
\end{enumerate}

\subsection{挑战冲刺训练}

\textbf{12.} （工业流程综合）工业上以软锰矿（主要成分为 $\text{MnO}_2$，含 $\text{Fe}_2\text{O}_3$、$\text{Al}_2\text{O}_3$、$\text{SiO}_2$ 等杂质）为原料制备 $\text{MnSO}_4\cdot\text{H}_2\text{O}$ 的流程如下：

软锰矿 $\xrightarrow{\text{稀H}_2\text{SO}_4 + \text{FeSO}_4}$ 酸浸 $\to$ 过滤 $\xrightarrow{\text{MnO}_2}$ 氧化 $\xrightarrow{\text{调pH}}$ 过滤 $\xrightarrow{\text{MnS}}$ 除杂 $\to$ 过滤 $\xrightarrow{\text{蒸发浓缩、冷却结晶}}$ $\text{MnSO}_4\cdot\text{H}_2\text{O}$

\begin{enumerate}
  \item 酸浸时 $\text{MnO}_2$ 被还原为 $\text{Mn}^{2+}$，离子方程式：\underline{\qquad\qquad}。
  \item 加入 $\text{MnO}_2$ 氧化时发生反应的离子方程式：\underline{\qquad\qquad}。
  \item 调 pH 的目的是\underline{\qquad\qquad}，控制 pH 的范围是\underline{\qquad\qquad}。
  \item 加入 $\text{MnS}$ 除杂时反应的离子方程式：\underline{\qquad\qquad}。
  \item 简述 $\text{MnSO}_4\cdot\text{H}_2\text{O}$ 的检验方法：\underline{\qquad\qquad}。
  \item 若 $1\ \text{t}$ 软锰矿中 $\text{MnO}_2$ 的质量分数为 87\%，理论上可制得 $\text{MnSO}_4\cdot\text{H}_2\text{O}$ 的质量为\underline{\qquad\qquad} $\text{kg}$。
\end{enumerate}

\textbf{13.} （化学反应原理综合）甲醇 $\text{CH}_3\text{OH}$ 可通过 $\text{CO}$ 和 $\text{H}_2$ 合成：
\[
\text{CO}(g) + 2\text{H}_2(g) \rightleftharpoons \text{CH}_3\text{OH}(g)\ \Delta H < 0
\]

\begin{enumerate}
  \item $T^\circ\text{C}$ 下，向 $2\ \text{L}$ 恒容容器中充入 $1\ \text{mol}\ \text{CO}$ 和 $2\ \text{mol}\ \text{H}_2$，平衡时 $\text{CO}$ 的转化率为 50\%，则该温度下的 $K =$\underline{\qquad\qquad}。
  \item 维持温度不变，再向容器中充入 $1\ \text{mol}\ \text{CO}$ 和 $2\ \text{mol}\ \text{H}_2$，$\text{CO}$ 的转化率\underline{\qquad\qquad}（填"增大""减小"或"不变"），新平衡时 $\text{CH}_3\text{OH}$ 的浓度为\underline{\qquad\qquad} $\text{mol/L}$。
  \item 利用甲醇燃料电池可实现化学能向电能的转化，写出酸性条件下该电池的负极反应式：\underline{\qquad\qquad}。
\end{enumerate}

\textbf{14.} （实验探究综合）某小组探究 $\text{Fe}^{2+}$ 和 $\text{Fe}^{3+}$ 的相互转化。

\begin{enumerate}
  \item 向 $\text{FeCl}_3$ 溶液中加入过量 $\text{Fe}$ 粉，现象：\underline{\qquad\qquad}，离子方程式：\underline{\qquad\qquad}。
  \item 取反应后上层清液，滴加 $\text{KSCN}$ 溶液，现象：\underline{\qquad\qquad}。
  \item 取反应后上层清液，滴加 $\text{H}_2\text{O}_2$ 溶液，现象：\underline{\qquad\qquad}，离子方程式：\underline{\qquad\qquad}。
  \item 设计实验证明 $\text{Fe}^{3+}$ 的氧化性比 $\text{I}_2$ 强：\underline{\qquad\qquad}。
\end{enumerate}

\subsection{限时训练}

\subsubsection{限时训练一（30 分钟）}

\begin{enumerate}
  \item 设 $N_A$ 为阿伏加德罗常数的值。下列说法正确的是（3 分）
  A. $0.1\ \text{mol}\ \text{H}_2\text{O}_2$ 含有的共价键数目为 $0.3N_A$
  B. $1\ \text{L}\ 0.1\ \text{mol/L}\ \text{Na}_2\text{SiO}_3$ 溶液中含有的 $\text{SiO}_3^{2-}$ 数目为 $0.1N_A$
  C. $2.4\ \text{g}\ \text{Mg}$ 在空气中充分燃烧，转移电子数目为 $0.1N_A$
  D. 标准状况下，$2.24\ \text{L}\ \text{C}_2\text{H}_5\text{OH}$ 含有的 $\text{C-H}$ 键数目为 $0.5N_A$
  \end{enumerate}

  \textbf{2.} 下列离子方程式正确的是（3 分）
  A. $\text{Na}$ 与 $\text{C}_2\text{H}_5\text{OH}$ 反应：$2\text{Na} + 2\text{C}_2\text{H}_5\text{OH} = 2\text{C}_2\text{H}_5\text{ONa} + \text{H}_2\uparrow$
  B. $\text{Fe(OH)}_2$ 与稀 $\text{HNO}_3$ 反应：$\text{Fe(OH)}_2 + 2\text{H}^+ = \text{Fe}^{2+} + 2\text{H}_2\text{O}$
  C. 向 $\text{Na}_2\text{S}_2\text{O}_3$ 溶液中加入稀 $\text{H}_2\text{SO}_4$：$2\text{S}_2\text{O}_3^{2-} + 4\text{H}^+ = 2\text{S}\downarrow + 2\text{SO}_2\uparrow + 2\text{H}_2\text{O}$
  D. $\text{NH}_4\text{HSO}_4$ 与等物质的量的 $\text{Ba(OH)}_2$ 反应：$2\text{H}^+ + \text{SO}_4^{2-} + \text{Ba}^{2+} + 2\text{OH}^- = \text{BaSO}_4\downarrow + 2\text{H}_2\text{O}$

\textbf{3.} 短周期元素 $W$、$X$、$Y$、$Z$ 的原子序数依次增大，$W$ 的单质为密度最小的气体，$X$ 的原子最外层电子数是其内层电子数的 3 倍，$Y$ 的原子半径在同周期中最大，$Z$ 的原子序数比 $Y$ 大 9。下列说法错误的是（3 分）
A. 原子半径：$Y > Z > X > W$
B. $W$ 与 $X$ 可形成两种常见化合物
C. $Y$ 的最高价氧化物对应水化物是强碱
D. $Z$ 的简单气态氢化物热稳定性比 $X$ 的强

\textbf{4.} 下列实验方案能达到实验目的的是（3 分）
A. 用 $\text{NaOH}$ 溶液除去 $\text{CO}_2$ 中的 $\text{SO}_2$
B. 用 $\text{Na}_2\text{CO}_3$ 溶液鉴别乙醇、乙酸和乙酸乙酯
C. 用 $\text{苯}$ 萃取溴水中的溴后，从分液漏斗上口放出
D. 用 $\text{FeCl}_3$ 溶液检验 $\text{FeSO}_4$ 溶液是否变质

\textbf{5.} 一种用于驱动潜艇的 $\text{Li}-\text{CO}_2$ 电池工作原理如图。下列说法正确的是（3 分）
A. 放电时 $\text{Li}$ 电极发生还原反应
B. 放电时 $\text{CO}_2$ 在正极得电子转化为 $\text{C}$ 和 $\text{Li}_2\text{CO}_3$
C. 放电时电解质中 $\text{Li}^+$ 向负极移动
D. 该电池的总反应为 $4\text{Li} + \text{CO}_2 = 2\text{Li}_2\text{O} + \text{C}$

\textbf{6.} （12 分）工业上以重晶石（$\text{BaSO}_4$）为原料制备 $\text{BaCl}_2$ 的流程如下：

重晶石 $\xrightarrow{\text{焦炭、煅烧}}$ 固体 $\xrightarrow{\text{HCl}}$ $\text{BaCl}_2$ 溶液 $\xrightarrow{\text{结晶}}$ $\text{BaCl}_2\cdot2\text{H}_2\text{O}$

\begin{enumerate}
  \item 煅烧时发生反应的化学方程式：\underline{\qquad\qquad}。
  \item 固体中主要成分是\underline{\qquad\qquad}。
  \item 用 $\text{HCl}$ 酸溶的离子方程式：\underline{\qquad\qquad}。
  \item 从 $\text{BaCl}_2$ 溶液得到 $\text{BaCl}_2\cdot2\text{H}_2\text{O}$ 的操作是\underline{\qquad\qquad}。
  \item $\text{BaCl}_2$ 溶液中混有少量 $\text{Fe}^{3+}$，可加入\underline{\qquad\qquad}除去。
\end{enumerate}

\textbf{7.} （12 分）已知 $A$、$B$、$C$、$D$ 均为有机物，$A$ 的产量是衡量国家石油化工水平的重要标志。转化关系如下：

$A \xrightarrow{\text{H}_2\text{O}} B \xrightarrow{\text{O}_2/\text{Cu}} C \xrightarrow{\text{O}_2/\text{催化剂}} D$\\
$B + D \xrightarrow{\text{浓H}_2\text{SO}_4,\Delta} E(\text{C}_4\text{H}_8\text{O}_2)$

\begin{enumerate}
  \item $A$ 的结构简式：\underline{\qquad\qquad}。
  \item $B \to C$ 的化学方程式：\underline{\qquad\qquad}。
  \item $B + D \to E$ 的化学方程式：\underline{\qquad\qquad}，反应类型：\underline{\qquad\qquad}。
  \item 鉴别 $B$ 和 $D$ 的常用试剂：\underline{\qquad\qquad}，现象：\underline{\qquad\qquad}。
\end{enumerate}
